\chapter{Le coût d'une page lue}\label{ch:cout}

La grille du chapitre précédent pose des questions d'accès, de catalogue, de droits et de
lecture. Il en manque une, que les détenteurs posent de plus en plus : combien coûte, en
énergie et en carbone, une lecture par la machine ? L'objection faite à ces projets par
les gardiens d'une archive mathématique (chapitre~\ref{ch:bourbaki}) tenait à l'éthique
autant qu'aux ressources. La part des ressources peut se mesurer. Le projet Grothendieck
l'a fait en novembre 2026, pour ses trois premiers mois (du 7 août au 7 novembre 2026) ;
la méthode est dans \emph{Mathématiques outillées}, et le détail dans l'article sur le
coût carbone du projet (\fichier{docs/carbon/article.tex}). Ce chapitre en tire ce qui
regarde les fonds.

\section{Un chiffre par page}

Pour tout le projet, transcriptions, lectures modernisées, preuves Lean et articles
compris, l'estimation centrale est d'environ 110~kg\,CO$_2$e, dans une fourchette de 15~kg
à 1,1~t. Presque tout vient de l'inférence du modèle ; l'intégration continue compte pour
1 à 4~kg. Divisée par les 9\,414 pages lues, l'inférence donne environ 12~g\,CO$_2$e par
page, de 2 à 120~g selon les hypothèses. C'est un majorant pour la seule transcription,
puisqu'une partie du calcul est allée au reste. Si l'on compte la fabrication des serveurs,
il faut multiplier par cinq à neuf.

Une page transcrite à la main par une ou un spécialiste porterait, à deux à cinq pages par
jour et deux cents jours de travail par an, de 2 à 12~kg\,CO$_2$e, si l'on y imputait toute
l'empreinte professionnelle de la personne (les 3,6~t par an et par membre de laboratoire
trouvées par l'outil GES~1point5, Mariette et al., 2022,
\url{https://doi.org/10.1088/2634-4505/ac84a4}). Le rythme est une hypothèse, que seuls les
éditeurs peuvent confirmer.

Les deux chiffres ne se divisent pas l'un par l'autre. Le chiffre humain impute des
bâtiments, des voyages et des achats qui servent à tout le travail d'un chercheur. Surtout,
les deux ne sont pas le même produit : une page lue par la machine est une première
lecture, que personne n'a encore vérifiée contre le fac-similé ; une édition est un texte
établi. La lecture par la machine prépare le travail humain, elle ne le remplace pas, et
son coût s'ajoute à celui d'une édition.

\section{D'une source à l'autre, le coût change}

Le coût d'une page dépend moins de la page que de la manière dont on la lit. Le modèle
relit à chaque tour tout ce qu'on lui a déjà donné : sur la fenêtre mesurée (du 24 septembre
au 7 novembre 2026), chaque appel au modèle principal relisait en moyenne environ 515\,000
jetons pour en écrire environ 680. La source décide de la taille de ce qu'on lui donne.

\begin{description}
  \item[La main.] Une dactylographie se lit en lots de quarante pages et en une passe ; un
        brouillon de Grothendieck demande des lots de vingt, des tuiles à 700~dpi et une
        seconde passe sur les agrandissements (chapitre~\ref{ch:montpellier}). Chaque image
        de plus, chaque passe de plus, est relue aux tours suivants.
  \item[La résolution.] On lit le texte à la plus basse résolution qui suffit (110~dpi
        pour les rédactions dactylographiées) et l'on garde les agrandissements pour les
        formules, les ajouts à la main et les lectures douteuses.
  \item[Ce qui est déjà lu.] Une page transcrite et commitée ne se relit pas pour une autre
        tâche : la lecture modernisée part de la transcription, pas de l'image. Une page
        écrite au dos d'un document sans rapport est sautée.
  \item[Le stockage.] Les fac-similés sont servis par le détenteur et lus en direct, non
        copiés sur un serveur du projet. La copie de travail locale ne sert qu'à
        transcrire. On évite ainsi un stockage en double, et c'est aussi ce que demandent
        les droits.
\end{description}

\begin{regle}
Avant de lancer les lots d'une source, on note dans son issue une estimation de ce que la
lecture coûtera : le nombre de pages, la main, la résolution nécessaire, le nombre de
passes. On la compare ensuite à ce que les journaux de l'assistant ont enregistré, avec la
date du relevé.
\end{regle}

\section{Choisir ce qu'on lit}

La machine a lu 9\,414 pages en trois mois, presque deux fois les 5\,129 pages que les
éditions faites à la main avaient établies dans ce fonds. Sans elle, ces pages n'auraient
pas été transcrites en trois mois. Le calcul n'a donc pas remplacé un travail qui se serait
fait autrement : il a servi à lire davantage. C'est l'effet rebond, et il interdit de
présenter un coût par page faible comme une économie. Le total croît avec ce qu'on décide
de lire.

La règle de sélection d'une source (chapitre~\ref{ch:premier-mois}) sert aussi à cela. Une
première lecture de tout un fonds a une valeur, celle d'un inventaire du contenu, qui aide
à choisir ce qu'il faudra éditer avec soin. La même lecture refaite pour une nouvelle mise
en page, ou une transcription relue sans raison, n'en a guère. Les pages déjà éditées à la
main n'ont pas à être relues par la machine ; on renvoie à l'édition.

\section{Les inconnues}

L'énergie par jeton d'un modèle commercial n'est pas publiée, ni l'emplacement des centres
de données ; les fourchettes viennent de là. L'entraînement des modèles est laissé de côté.
L'eau n'est pas estimée : les deux chiffres publiés par réponse diffèrent de plus d'un
facteur cent (0,26~mL pour la requête médiane de Gemini chez Google, 45~mL pour une réponse
de 400 jetons chez Mistral AI), et rien ne permet de choisir. Moins de la moitié de la
période a été mesurée directement, parce que l'outil effaçait ses journaux après quatorze
jours ; la conservation a été allongée le 7 novembre 2026.

\section{Répondre à l'objection}

Un chiffre ne répond qu'à la part environnementale d'une objection. Il ne dit rien de la
part éthique : à qui profite l'outil, qui a consenti à ce que les archives soient lues par
lui, ce que devient le texte. Un détenteur peut refuser pour ces raisons-là, et son refus
vaut, quel que soit le nombre de grammes par page. Le projet peut seulement apporter à la
discussion un compte daté, sa méthode, et la liste des gestes qui le font baisser.

\begin{exemple}
Trois mois de lecture d'un fonds de 9\,414 pages ont coûté de l'ordre de quelques pour cent
de l'empreinte professionnelle annuelle d'un chercheur, et plusieurs fois plus avec le
matériel. La plus grande part tient à la longueur des contextes relus, et non au nombre de
pages. C'est sur elle qu'on agit d'abord.
\end{exemple}

\begin{outils}
Les journaux de l'assistant et \cmd{gh api} pour mesurer (le détail est dans
\emph{Mathématiques outillées}) ; l'issue de la source pour l'estimation et le relevé ;
le fac-similé servi par le détenteur plutôt qu'une copie hébergée.
\end{outils}
